\BourbakiExerciseGroup

\begin{enumerate}
\def\labelenumi{\arabic{enumi})}
\item
  In a Hausdorff topological ring A, the commutant of every subset of A (and in particular the centre of A) is closed, and so is the left (resp. right) annihilator of every subset of A.
\item
  \begin{enumerate}
  \def\labelenumii{\alph{enumii})}
  \tightlist
  \item
    Let \(\mathfrak{a}\) be a discrete left ideal in a topological ring A. Show that, for every \(x \in \mathfrak{a}\) , the left annihilator of \(x\) in A is open. Hence show that if A is a non-discrete ring with no zero-divisors, it contains no discrete (left or right) ideal other than \(\{\mathbf{o}\}\) .
  \end{enumerate}
\end{enumerate}

\begin{enumerate}
\def\labelenumi{\alph{enumi})}
\setcounter{enumi}{1}
\tightlist
\item
  Let A be a locally connected topological ring, and let \(\mathfrak{a}\) be a totally disconnected left ideal of A. Show that, for each \(x \in \mathfrak{a}\) , the left annihilator of \(x\) in A is open. Hence show that if A has no zero-divisors it contains no totally disconnected (left or right) ideal other than \(\{0\}\) .
\end{enumerate}

\begin{enumerate}
\def\labelenumi{\arabic{enumi})}
\setcounter{enumi}{2}
\item
  The connected component of o in a topological ring A is a two-sided ideal. Deduce that if A is quasi-simple and is not connected (in particular, if A is a non-connected topological division ring) then A is totally disconnected.
\item
  Let \((x_{i})\) be a summable family in a Hausdorff topological ring A, and let \(s\) be its sum. Then for each \(a \in A\) the family \((ax_{i})\) {[}resp. \((x_{i}a)]\) is summable and its sum is as (resp. sa). If \((x_{\lambda})\) and \((y_{\mu})\) are two summable families in A such that the family \((x_{\lambda}y_{\mu})\) is summable, then \(\sum_{\lambda,\mu} x_{\lambda} y_{\mu} = \left( \sum_{\lambda} x_{\lambda} \right) \left( \sum_{\mu} y_{\mu} \right)\) . If moreover A is complete and if one of the \(x_{\lambda}\) is a unit, then the summability of the family \((x_{\lambda}y_{\mu})\) implies the summability of \((y_{\mu})\) .
\item
  A rectangle is a subset of \(N \times N\) which is the product of two intervals of N. Given a finite subset E of \(N \times N\) , let \(\varphi(E)\) be the smallest number of mutually disjoint rectangles whose union is E. Let \((E_n)\) be an increasing sequence of finite subsets of \(N \times N\) which cover \(N \times N\) and are such that the sequence \((\varphi(E_n))\) is bounded. If \((x_n), (y_n)\) are two convergent series whose terms belong to a Hausdorff topological ring A, then we have
\end{enumerate}

\[
\lim _ {n \rightarrow \infty} \sum_ {(h, k) \in \mathrm{E} _ {n}} x _ {h} y _ {k} = \left(\underset {n = 0} {\overset {\infty} {\sum}} x _ {n}\right)\left(\underset {n = 0} {\overset {\infty} {\sum}} y _ {n}\right).
\]

Construct an example of a sequence \((\mathbf{E}_{n})\) satisfying the conditions above and such that \(E_{n+1} - E_{n}\) contains only one element for each n.

\begin{enumerate}
\def\labelenumi{\arabic{enumi})}
\setcounter{enumi}{5}
\item
  Let A be a topological ring with an identity element and let \(a, b\) be two two-sided ideals of A such that \(a + b = A\) and \(a \cap b = \{0\}\) . Show that the topological ring A is isomorphic to the product of the topological rings a and b.
\item
  Let A be a topological ring with no identity element, and let B be the ring obtained by adjoining an identity element to A show that the topology on B which is the product of the topology of A and the discrete topology on Z is compatible with the ring structure of B, and induces the given topology on A (a two-sided ideal of B).
\item
  Let \(\mathbf{A}\) be a topological ring with an identity element.
\end{enumerate}

\begin{enumerate}
\def\labelenumi{\alph{enumi})}
\item
  The set A* of units of A is open in A if and only if there is a neighbourhood V of 1 all of whose elements are units in A.
\item
  If A* is open in A then every maximal (left or right) ideal of A is closed, and the radical of A is closed. If A has no closed left ideals other than \(\{o\}\) and A, then A is a division ring.
\item
  Let A be the subring of the topological field Q consisting of all rational numbers \(k/2^{n}\) ( \(k \in Z, n \in N\) ). Show that A contains no closed ideal other than \(\{0\}\) and A, but is not a field.
\end{enumerate}

\begin{enumerate}
\def\labelenumi{\arabic{enumi})}
\setcounter{enumi}{8}
\tightlist
\item
  \begin{enumerate}
  \def\labelenumii{\alph{enumii})}
  \tightlist
  \item
    An element \(x\) (resp. an ideal \(\mathfrak{a}\) ) of a topological ring \(\mathbf{A}\) is said to be topologically nilpotent if the sequence \((x^n)_{n\geqslant 1}\) converges to o {[}resp. if the filter base formed by the \(\mathfrak{a}^n\) ( \(n\geqslant 1\) ) converges to o). Every element of a topologically nilpotent ideal is topologically nilpotent.
  \end{enumerate}
\end{enumerate}

\begin{enumerate}
\def\labelenumi{\alph{enumi})}
\setcounter{enumi}{1}
\item
  Suppose that A has an identity element and that \(A^{*}\) is open in A. Show that if x is any topologically nilpotent element of A, then 1 - x is a unit (observe that \(1 - x^{n}\) is a unit for sufficiently large n). Hence show that if all the elements of a (left or right) ideal a are topologically nilpotent, then a is contained in the radical of A.
\item
  Suppose that A is complete and has an identity element, and that there exists a fundamental system of neighbourhoods of o which are subgroups of the additive group of A. Show that if x is topologically nilpotent, then 1 - x is a unit in A.
\end{enumerate}

\(d)\) Under the hypotheses of \(c\) ), show that if \(y \in A\) is topologically nilpotent, then the equation \(x^{2} + x = y\) has a root \(x \in A\) , and that \(x\) is topologically nilpotent.

¶ 10) a) Let B be a ring with an identity element, let E be a free B-module \(\mathbf{B}_{d}^{(I)}\) , and let A be the ring \(\mathcal{L}(\mathbf{E})\) of endomorphisms of E. For every finite subset F of E, let \(V_{F}\) be the set of all \(u \in A\) such that \(u(x) = 0\) for all \(x \in F\) . Show that the \(V_{F}\) from a fundamental system of neighbourhoods of o for a Hausdorff topology on A, which is compatible with the ring structure of A, and with respect to which A is totally disconnected.

\begin{enumerate}
\def\labelenumi{\alph{enumi})}
\setcounter{enumi}{1}
\tightlist
\item
  Take I = N and take B to be a commutative ring which has no zero-divisors but whose radical R is not \(\{0\}\) . Show that the radical of A is not closed. {[}Let \((e_{n})\) be the canonical basis of E. Remark on the one hand that each \(u \in A\) , such that \(u(e_{n}) = 0\) for all but a finite number of indices and such that \(u(E) \subset \Re \cdot E\) , belongs to the radical of A; consider on the other hand the element \(u_{0}\) of A such that \(u_{0}(e_{n}) = e_{n} + te_{n+1}\) with \(t \neq 0\) in R, for each n{]}.
\end{enumerate}

¶ 11) a) A topological ring with an identity element is said to be a Gelfand ring if A* is open in A and if the topology induced on A* by the topology of A is compatible with the multiplicative group structure of A*. Every Hausdorff topological division ring is a Gelfand ring {[}cf.~Exercise 20 e){]}.

\begin{enumerate}
\def\labelenumi{\alph{enumi})}
\setcounter{enumi}{1}
\tightlist
\item
  Show that if A is a Gelfand ring then so is every matrix ring \(\mathbf{M}_{n}(\mathbf{A})\) endowed with the product topology (on \(A^{n^{2}}\) ). (Proof by induction on n).
\end{enumerate}

\begin{enumerate}
\def\labelenumi{\arabic{enumi})}
\setcounter{enumi}{11}
\tightlist
\item
  A subset M of a topological ring is said to be right bounded (resp. left bounded) if, for each neighbourhood U of o in A, there is a neighbourhood V of o such that VM⊂U (resp. MV⊂U); M is bounded if it is both left and right bounded. The topology of A is said to be locally bounded (and A is said to be a locally bounded topological ring) if there is a bounded neighbourhood of o in A.
\end{enumerate}

\begin{enumerate}
\def\labelenumi{\alph{enumi})}
\item
  Every topological ring which has a fundamental system of neighbourhoods of o consisting of right ideals is right bounded. Show that if I is infinite and B is a division ring in Exercise 10 a), then the ring A is left bounded, but there is no right bounded neighbourhood of o in A.
\item
  If a topological ring A is right bounded (resp. bounded) and has a fundamental system of neighbourhoods of o which are subgroups of the additive group of A, then it has a fundamental system of neighbourhoods of o which are right (resp. two-sided) ideals.
\item
  Every finite union of bounded sets is bounded; the closure of a bounded set is bounded; if M and N are bounded, then so are \(M + N\) and MN.
\item
  Every precompact set in a topological ring \(\mathbf{A}\) is bounded.
\item
  If M and N are two bounded subsets of A, show that the mapping \((x, y) \to xy\) elements of \(M \times N\) into A is uniformly continuous.
\item
  Let A be a Hausdorff topological ring which has an identity element. Show that if \((x_{n})\) is a sequence of units which tends to o in A, then the set of elements \(x_{n}^{-1}\) is neither left nor right bounded.
\item
  If A is a bounded ring with an identity element, show that the mapping \(x \rightarrow x^{-1}\) is uniformly continuous on the set A* of units of A. h) Show that if A is a complete Hausdorff bounded ring with an identity element, then the set A* of units and the radical of A are closed {[}use g){]}.
\item
  A subset of a product \(\prod_{t} A_{t}\) of topological rings is bounded if and only if all its projections are bounded.
\item
  Every subring of a locally bounded ring and every product of locally bounded rings are locally bounded rings.
\item
  The completion of a locally bounded Hausdorff ring is locally bounded.
\end{enumerate}

\begin{enumerate}
\def\labelenumi{\arabic{enumi})}
\setcounter{enumi}{12}
\tightlist
\item
  Let A be a bounded topological ring (Exercise 12) which has an identity element and is such that the set A* of units of A is open in A. Show that the radical of A is open. Hence if A is without radical, it is discrete; in particular, a bounded Hausdorff topological division ring is discrete. A compact ring A without radical, which has an identity element and in which A* is open, is finite; in particular, a compact topological division ring is finite.
\end{enumerate}

¶ 14) Let A be a compact ring with an identity element, totally disconnected (*) and without radical. Then A is isomorphic to the product of a family of finite simple rings. {[}Show first that A has a fundamental system of neighbourhoods of o which are two-sided ideals, using Exercise 12 b) of § 6 and Proposition 14 of § 4. Hence show that there is a maximal set \(\Phi\) of two-sided ideals of A such that (i) every two-sided ideal \(\mathfrak{M} \in \Phi\) is a maximal open two-sided ideal of A; (ii) no ideal \(\mathfrak{M} \in \Phi\) contains a finite intersection of ideals of \(\Phi\) other than \(\mathfrak{M}\). Show then that A is isomorphic to the product of the rings \(A/\mathfrak{M}\) as \(\mathfrak{M}\) runs through \(\Phi\); prove first that the intersection of all the \(\mathfrak{M} \in \Phi\) consists only of the zero element of A by using Exercise 9 b){]}.
(*) It can be shown that a compact ring with an identity element is necessarily totally disconnected.

\begin{enumerate}
\def\labelenumi{\arabic{enumi})}
\setcounter{enumi}{14}
\item
  Let A be a totally disconnected compact ring with an identity element. Show that the radical \(\Re\) of A is topologically nilpotent {[}Exercise 9 a); use Exercise 12 b) of § 6 and Proposition 14 of § 4{]}. Hence show that \(\Re\) is the set of all \(x \in A\) such that \(ax\) is topologically nilpotent for all \(a \in A\) {[}cf.~Exercise 9 c){]}.
\item
  Show that on a ring A (resp. a division ring K), the least upper bound \(\mathcal{C}\) of a family \((\mathcal{T}_i)\) of topologies compatible with the ring structure of A (resp. the division ring structure of K) is compatible with this structure. A set M which is left (resp. right) bounded for each of the topologies \(\mathcal{T}_i\) is also left (resp. right) bounded for \(\mathcal{C}\) .
\item
  If K is a Hausdorff topological division ring and is not discrete, then neither of the multiplicative uniformities on K is comparable with the uniformity induced on \(K^{*}\) by the additive uniformity of K (cf.~Exercise 13).
\item
  Let K be a Hausdorff topological division ring, let A be a closed subset of K and let B be a compact subset of K such that \(o \notin B\) . Show that AB and BA are closed in K (cf.~§ 4, no. 1, Proposition 1, Corollary 1). * In the field R of real numbers, give an example where A is closed, and B is compact, \(o \in B\) and AB is not closed. *
\item
  \begin{enumerate}
  \def\labelenumii{\alph{enumii})}
  \tightlist
  \item
    Let K be a division ring endowed with a Hausdorff topology \(\mathcal{G}\) compatible with the ring structure of K. Show that there is a neighbourhood V of o in K such that \((\mathbb{CV}) \cup (\mathbb{CV})^{-1} = \mathbf{K}^*\) {[}this is equivalent to \(\mathbf{V} \cap (\mathbf{V} \cap \mathbf{K}^*)^{-1} = \emptyset\) {]}.
  \end{enumerate}
\end{enumerate}

\begin{enumerate}
\def\labelenumi{\alph{enumi})}
\setcounter{enumi}{1}
\tightlist
\item
  Suppose moreover that \(\mathfrak{C}\) is not discrete. Show that if \(\mathbf{U}\) is any neighbourhood of \(o\) in \(\mathbf{K}\) and if \(\mathbf{U}^*\) denotes \(\mathbf{U} \cap \mathbf{K}^*\) , then
\end{enumerate}

\[
\mathbf {U} ^ {*} (\mathbf {U} ^ {*}) ^ {- 1} = \mathbf {K} ^ {*}.
\]

¶ 20) a) Let K be a division ring endowed with a non-discrete Hausdorff topology \(\mathfrak{C}\) compatible with the ring structure of K. A subset M of K is right (resp. left) bounded if and only if, for each neighbourhood U of o there exists \(a \in \mathbf{K}^*\) such that \(a\mathbf{M} \subset \mathbf{U}\) (resp. Ma \(\subset\) U); M is then right (resp. left) bounded for every Hausdorff topology \(\mathfrak{C}'\) which is coarser than \(\mathfrak{C}\) and is compatible with the ring structure of K.

\begin{enumerate}
\def\labelenumi{\alph{enumi})}
\setcounter{enumi}{1}
\item
  If \(\mathfrak{G}\) is locally bounded (Exercise 12) and if U is a bounded neighbourhood of o with respect to \(\mathfrak{G}\) , then the sets \(x\mathrm{U}\) (resp. \(\mathrm{U}x\) ) form a fundamental system of neighbourhoods of o (with respect to \(\mathfrak{G}\) ) as x runs through \(\mathbf{K}^*\) . Every point of K has a countable fundamental system of neighbourhoods of o with respect to \(\mathfrak{G}\) if and only if there is a sequence \((a_n)\) of points of K which has o as a cluster point.
\item
  The least upper bound of any finite family of locally bounded topologies on K (compatible with the ring structure of K) is locally bounded. Conversely, if the least upper bound \(\mathcal{G}\) of a family \((\mathfrak{T}_i)\) of locally bounded topologies on K is locally bounded, then \(\mathcal{G}\) is the least upper bound of a finite subfamily of \((\mathfrak{T}_i)\) . (Remark that if U is a neighbourhood of o which is bounded with respect to \(\mathcal{G}\) , then there is a finite system of indices \((\iota_k)\) and for each \(k\) a neighbourhood \(\mathrm{V}_k\) of o, bounded with respect to \(\mathcal{G}_{\iota_k}\) , such that \(\bigcap \mathrm{V}_k \subset \mathrm{U}\) ; on the other hand, there exists \(a_k \in \mathrm{K}^*\) such that \(\mathrm{U} \subset a_k \mathrm{V}_{k'}\) .)
\item
  A subset F of a division ring K is said to be a quasi-ring if: (i) \(F \neq K\) , \(o \in F\) , \(1 \in F\) , \(-F = F\) , \(FF \subset F\) ; (ii) there is an element \(c \in F^{*} = F \cap K^{*}\) such that \(c(F + F) \subset F\) ; and (iii) for each \(x \in F\) , there exists \(y \in F\) such that \(yF \subset Fx\) and \(Fy \subset xF\) . Show that if G is a non-discrete Hausdorff topology which is locally bounded and compatible with the ring structure of K, then every symmetric bounded neighbourhood U of o (with respect to G), such that \(UU \subset U\) and \(1 \in U\) , is a quasi-ring, and that \(U^{*}(U^{*})^{-1} = K^{*}\) {[}Exercise 19 b){]}. If V is any symmetric bounded neighbourhood of o with respect to G, then the set U of elements \(x \in K\) such that \(xV \subset V\) (or \(Vx \subset V\) ) is a symmetric bounded neighbourhood of o with respect to G such that \(UU \subset U\) and \(1 \in U\) , and hence is a quasi-ring.
\item
  Conversely, let F be a quasi-ring in K such that \(\mathrm{F}^{*}(\mathrm{F}^{*})^{-1} = \mathrm{K}^{*}\) . Show that there is a Hausdorff topology G which is not discrete and is compatible with the ring structure of K, such that F is a bounded neighbourhood of o with respect to G. In particular we may take F to be any subring A of K such that K is the division ring of right fractions of A. In particular, if K = Q and F = Z, the corresponding topology G on K is not compatible with the field structure of K.
\end{enumerate}

¶ 21) a) Let \(\mathfrak{C}\) be a locally bounded Hausdorff topology on a division ring K, with respect to which K is connected. Show that there are non-zero topologically nilpotent elements {[}Exercise 9 a){]} in K. (If U is a symmetric bounded neighbourhood of o, and t is a non-zero element of K such that \(\mathrm{U}t + t\mathrm{U} \subset \mathrm{U}\) , show that K is the union of the sets \(\mathrm{U}t^{-n}\) , by using Proposition 6 of § 2, no. 2).

\begin{enumerate}
\def\labelenumi{\alph{enumi})}
\setcounter{enumi}{1}
\tightlist
\item
  Let \(\mathfrak{G}\) be a locally bounded Hausdorff topology on K, with respect to which K has non-zero topologically nilpotent elements. Show that there is a neighbourhood U of o such that, for each neighbourhood V of o, there is an integer \(n_0\) such that \(\mathrm{U}^n\subset \mathrm{V}\) whenever \(n\geqslant n_0\) ; this implies that all the elements of U are topologically nilpotent and that every point of K has a countable fundamental system of neighbourhoods with respect to \(\mathfrak{G}\) . (Let \(t\neq 0\) be topologically nilpotent and let W be a bounded neighbourhood of o; take U to be a bounded neighbourhood of o such that \(\mathrm{UW}\subset \mathrm{Wt}\) .) The set T of topologically nilpotent elements of K is therefore a neighbourhood of o with respect to \(\mathfrak{G}\) .
\end{enumerate}

¶ 22) In a division ring K endowed with a Hausdorff topology \(\mathfrak{G}\) compatible with the ring structure of K, a set R containing o is said to be retrobounded if \((\mathbb{C}R)^{-1}\) is bounded. \(\mathfrak{G}\) is said to be locally retrobounded if there is a fundamental system of retrobounded neighbourhoods of o in K with respect to \(\mathfrak{G}\) .

\begin{enumerate}
\def\labelenumi{\alph{enumi})}
\item
  A locally retrobounded topology T is locally bounded, and is a minimal element in the set of all Hausdorff topologies compatible with the ring structure of K; also C is compatible with the division ring structure of K {[}use Exercise 19 a){]}.
\item
  Suppose from now on, in this exercise, that \(\mathcal{G}\) is a locally retrobounded topology on K. Show that, for every neighbourhood V of o in K, \(x \to x^{-1}\) is uniformly continuous on \(\mathbb{C}\mathrm{V}\) . Hence show that the completion \(\hat{\mathbf{K}}\) of K is a topological division ring whose topology is locally retrobounded.
\item
  Every neighbourhood of o in K (with respect to \(\mathfrak{G}\) ) contains a neighbourhood V of o such that \(\mathbf{E}(\mathbf{V}) = (\mathbb{C}\mathbf{V}) \cap (\mathbb{C}\mathbf{V})^{-1}\) is not empty. Show that the three uniformities on \(\mathbf{E}(\mathbf{V})\) , induced by the additive uniformity and the two multiplicative uniformities on K, coincide. Deduce that, if K is complete, then \(\mathbf{K}^*\) is complete with respect to either multiplicative uniformity.
\item
  Show that if \(x\) is an element of \(\mathbf{K}^*\) then exactly one of the following three statements is true: (i) \(x\) is topologically nilpotent; (ii) \(x^{-1}\) is topologically nilpotent; (iii) the sequence \((x^n)_{n\in \mathbb{Z}}\) is bounded. {[}By considering a bounded neighbourhood \(U\) of \(o\) such that \(UU\subset U\) (Exercise 20 d)), show that if o is a cluster point of the sequence \((x^n)_{n\geqslant 0}\) , then \(x\) is topologically nilpotent.{]} If K contains non-zero topologically nilpotent elements, show that the set B, consisting of o and all \(x\in \mathbf{K}^*\) such that \(x^{-1}\) is not topologically nilpotent, is a bounded neighbourhood of o in K which is invariant under all inner automorphisms {[}use Exercise 21 b){]}; B then contains all bounded neighbourhoods U of o such that \(\mathrm{UU}\subset \mathrm{U}\) ; the set of such neighbourhoods which are invariant under all inner automorphisms has a greatest element, which coincides with B if the commutator subgroup of \(\mathbf{K}^*\) is bounded, in particular if K is a field.
\item
  If U is a bounded symmetric neighbourhood of o in K such that \(UU \subset U\) , then there exists \(b \in U^{*}\) such that \(\mathbf{K}^{*} = \mathbf{U}^{*} \cup ((\mathbf{U}^{*})^{-1} b)\) . Let \(a \in U^{*}\) be such that \(Ua \subset bU\) , and suppose that the sequence \((a^{-n})_{n \geqslant 0}\) is bounded. Show that the set V of all \(x \in K\) such that xU is contained in the union of the sets \(Ua^{-n}\) is a bounded neighbourhood of o such that \(VV \subset V\) and \(\mathbf{K}^{*} = \mathbf{V}^{*} \cup (\mathbf{V}^{*})^{-1}\) .
\item
  Suppose that K contains no non-zero topologically nilpotent elements. Show that there is a bounded neighbourhood of o in K which is a subring A of K such that \(\mathbf{K}^{*} = \mathbf{A}^{*} \cup (\mathbf{A}^{*})^{-1}\) . {[}Begin with a bounded neighbourhood V of o such that \(\mathbf{K}^{*} = \mathbf{V}^{*} \cup (\mathbf{V}^{*})^{-1}\) (see e){]}; if \(c \in V^{*}\) is such that \(c(V + V) \subset V\) , observe that the ring A generated by V is contained in the union of the sets \(Vc^{-n}\) ( \(n \geqslant 0\) ).
\end{enumerate}

¶ 23) If p is any prime number and x is any non-zero rational number, let \(v_{p}(x)\) denote the exponent of p in the decomposition of x into prime factors. The mapping \(v_{p}\) of \(Q^{*}\) into Z is called the p-adic valuation on Q. The set of all \(x \in Q\) such that either x = 0 or \(v_{p}(x) \geqslant m\) (for some \(m \in Z\) ) is the fractional ideal ( \(p^{m}\) ) of Q.

\begin{enumerate}
\def\labelenumi{\alph{enumi})}
\tightlist
\item
  Show that the ideals \((p^{m})\) ( \(m \in Z\) ) form a fundamental system of neighbourhoods of o in Q for a locally retrobounded topology (Exercise 22) \(C_{p}\) on Q, called the p-adic topology. The completion \(Q_{p}\) of Q in this topology is a field called the p-adic field, whose elements are called p-adic numbers. The closure \(Z_{p}\) of Z in \(Q_{p}\) is a compact open principal ideal domain in \(Q_{p}\) , in which \(p = pZ_{p}\) is the only non-zero prime ideal; \(Z_{p}/p\) is isomorphic to the prime field \(F_{p} = Z/(p)\) , and more generally the additive quotient group \(p^{m}/p^{n}\) is isomorphic to
\end{enumerate}

\[
\mathbf {Z} / (p ^ {n - m}) \text {   if   } m <   n.
\]

The elements of \(\mathbf{Z}_p\) are called \(p\) -adic integers.

\begin{enumerate}
\def\labelenumi{\alph{enumi})}
\setcounter{enumi}{1}
\item
  Show that every continuous homomorphism of the additive group \(Q_{p}\) into itself is of the form \(x \to ax\) , where \(a \in Q_{p}\) {[}if f is such a homomorphism, show that \(f(rx) = rf(x)\) for all \(r \in Q\) , and then pass to the limit in \(Q_{p}\) {]}.
\item
  Show that every compact subgroup \(G \neq \{o\}\) of the additive group \(Q_{p}\) is identical with one of the \(p^{n} (n \in Z)\) , and that there is no non-compact closed subgroup of the additive group \(Q_{p}\) , except \(Q_{p}\) itself. {[}Let m be the largest integer such that \(G \subset p^{m}\) ; by considering the quotient group \((G + p^{n})/p^{n}\) for n \textgreater{} m, show that \(G + p^{n} = p^{m}\) , and then use formula (1) of §3, no. 1{]}.
\end{enumerate}

¶ 24) a) Show that every subgroup of the multiplicative group \(\mathbf{Q}_p^*\) of the \(p\) -adic field \(\mathbf{Q}_p\) is isomorphic to the product of a subgroup of the multiplicative group \(\mathbf{U}\) of units of \(\mathbf{Z}_p\) and a discrete additive subgroup isomorphic to \(\mathbf{Z}\) or \(\{o\}\) .

\begin{enumerate}
\def\labelenumi{\alph{enumi})}
\setcounter{enumi}{1}
\item
  Show that the compact subgroups of the subgroup \(V = 1 + p\) of U are the groups \(1 + p^{n}\) {[}same reasoning as in Exercise 23 c){]}.
\item
  Show that, for each \(a \in \mathbf{U}\) , the sequence \((a^{p^n})_{n \in \mathbb{N}}\) tends to a limit \(\alpha \equiv a \pmod{\mathfrak{p}}\) and that \(\alpha^p = \alpha\) {[}show that \(a^{p^n} \equiv a^{p^{n-1}} \pmod{\mathfrak{p}^n}\) by induction on \(n\) {]}. Show that all the roots of the polynomial \(\mathbf{X}^{p-1} - 1\) (in an algebraically closed extension of \(\mathbf{Q}_p\) ) belong to \(\mathbf{Q}_p\) and are pairwise incongruent mod \(\mathfrak{p}\) {[}apply what precedes to the roots of the congruence \(x^{p-1} - 1 \equiv 0 \pmod{p}\) {]}. If \(p > 2\) and if \(d\) is the highest common factor of \(n\) and \(p - 1\) , the polynomial \(\mathbf{X}^n - 1\) has exactly \(d\) roots in \(\mathbf{Q}_p\) , which are roots of \(\mathbf{X}^d - 1\) (same method: take \(a\) to be a root of \(\mathbf{X}^n - 1\) ).
\end{enumerate}

Hence show that, if \(p > 2\) , every compact subgroup of the multiplicative group \(\mathbf{Q}_p^*\) is the direct product of a finite subgroup of \(\mathbf{U}\) {[}consisting of the \((p - 1)\) th roots of unity{]} and a subgroup of the form \(1 + \mathfrak{p}^n\) {[}use \(b\) {]}. How are these results to be modified for \(p = 2\) ? (Make the group \(1 + \mathfrak{p}^2\) play the part previously played by \(1 + \mathfrak{p}\) ).

\begin{enumerate}
\def\labelenumi{(\arabic{enumi})}
\setcounter{enumi}{24}
\tightlist
\item
  \begin{enumerate}
  \def\labelenumii{\alph{enumii})}
  \tightlist
  \item
    Let \(a \neq o\) be a \(p\) -adic number. Then the mapping \(n \to a^n\) is continuous on \(\mathbf{Z}\) (as a subspace of \(\mathbf{Q}_p\) ) if and only if \(a \in 1 + \mathfrak{p}\) . If this condition is satisfied, show that \(n \to a^n\) is uniformly continuous on \(\mathbf{Z}\) and extends to a continuous homomorphism of \(\mathbf{Z}_p\) into U, denoted by \(x \to a^x\) , which is injective if \(a \neq 1\) . If \(p > 2\) and if \(m\) is the largest integer such that \(a \in 1 + \mathfrak{p}^m\) , then \(x \to a^x\) is an isomorphism of \(\mathbf{Z}_p\) onto \(1 + \mathfrak{p}^m\) . How must this statement be modified when \(p = 2\) ? b) Show that every continuous homomorphism of \(\mathbf{Z}_p\) into U is of the form \(x \to a^x\) for some \(a \in 1 + \mathfrak{p}\) {[}if \(f\) is such a homomorphism and \(f(1) = a\) , then \(f(n) = a^n\) for \(n \in \mathbf{Z}]\) .
  \end{enumerate}
\end{enumerate}

\begin{enumerate}
\def\labelenumi{\alph{enumi})}
\setcounter{enumi}{2}
\item
  Show that the mapping \((x, y) \to x^y\) is continuous on \((1 + \mathfrak{p}) \times \mathbf{Z}_p\) . If \(p > 2\) and \(b \in \mathbf{Z}_p\) , and if \(m\) is the largest integer such that \(b \in \mathfrak{p}^m\) , then the mapping \(x \to x^b\) is an isomorphism of the multiplicative group \(1 + \mathfrak{p}\) onto the multiplicative subgroup \(1 + \mathfrak{p}^{m+1}\) {[}use Exercise 24 b){]}. How must this statement be modified when \(p = 2\) ?
\item
  If \(n\) is an integer prime to both \(p - 1\) and \(p\) , show that the mapping \(x \to x^n\) is an automorphism of the multiplicative group U {[}use Exercise 24 c){]}.
\end{enumerate}

\begin{enumerate}
\def\labelenumi{\arabic{enumi})}
\setcounter{enumi}{25}
\tightlist
\item
  \begin{enumerate}
  \def\labelenumii{\alph{enumii})}
  \tightlist
  \item
    Let \(p, q\) be two distinct prime numbers and let \(\mathfrak{C}\) be the topology on \(\mathbf{Q}\) which is the least upper bound of the topologies \(\mathfrak{C}_p\) and \(\mathfrak{C}_q\) . Then \(\mathfrak{C}\) is locally bounded {[}Exercise 20 c){]} and is compatible with the field structure of \(\mathbf{Q}\) . With respect to this topology neither of the two sequences \(((p/q)^n)_{n\geqslant 0}, ((q/p)^n)_{n\geqslant 0}\) is bounded; the sequence \((1/p^n)_{n\geqslant 0}\) is not bounded, but the sequence \((p^n)_{n\geqslant 0}\) , which is bounded, does not tend to o {[}cf.~Exercise 22 d){]}. Show that the completion of \(\mathbf{Q}\) with respect to the topology \(\mathfrak{C}\) is isomorphic to the product of the topological fields \(\mathbf{Q}_p\) and \(\mathbf{Q}_q\) .
  \end{enumerate}
\end{enumerate}

\begin{enumerate}
\def\labelenumi{\alph{enumi})}
\setcounter{enumi}{1}
\tightlist
\item
  If P is the set of all prime numbers then the least upper bound \(G_{0}\) of the topologies \(G_{p}\) for \(p \in P\) is compatible with the field structure of Q but is not locally bounded {[}Exercise 20 c){]}. What is the completion of Q with respect to the topology \(G_{0}\) ?
\end{enumerate}

